\section{Design Details}
\label{sec:method}

\subsection{Problem Definition}

Given a protocol standard $S$ and an implementation $P$, \tool{} seeks to
identify whether $P$ satisfies a set of normative field-level requirements in
$S$. A requirement is represented as a rule
$r=\langle C, f, A, M, \Pi\rangle$, where $C$ is the applicability condition,
$f$ is the protocol field, $A$ is the expected action or constraint, $M$ is the
normative modality such as \emph{MUST} or \emph{SHOULD}, and $\Pi$ records
source provenance. The checker returns a tuple
$\langle y,E,U,T\rangle$: label $y\in\{\textsc{Consistent},
\textsc{Inconsistent},\textsc{Insufficient}\}$, evidence $E$, uncertainty notes
$U$, and optional validation-test artifact $T$.

\tool{} does not assume that the implementation uses names similar to the
standard. It also does not require a pre-existing formal model. Instead, it
builds a structured intermediate representation from the standard and aligns
that representation to code using a combination of semantic and syntactic
evidence.

\subsection{Field-Space Construction}
\label{sec:fieldspace}

Directly extracting rules from open-ended prose is unstable because standards
mix design motivation, compatibility advice, examples, and normative
requirements. \tool{} therefore starts by constructing a \emph{field space}: a
set of protocol objects that subsequent rules may reference. Field candidates
are collected from message formats, extension structures, parameter tables, and
explicit field-definition paragraphs. For each field, \tool{} records its
canonical name, aliases, parent message or structure, type, allowed values when
available, and source locations.

The field space serves two purposes. First, it gives the extractor a closed set
of referents, reducing the chance that descriptive concepts such as
``compatibility mode'' are mistaken for checkable fields. Second, it preserves
cross-document identity: the same field may appear as a full path, short name,
structure member, or contextual pronoun. During rule extraction, all references
must resolve to one or more field-space entries. Ambiguous references are kept
as uncertainty rather than silently collapsed.

\textbf{Field identity and validation.}
\tool{} treats two mentions as the same field only when their parent context
and semantic role are compatible, preventing generic names such as
\emph{length}, \emph{type}, or \emph{flags} from being merged across unrelated
messages. Before rule extraction, it also validates the field space by marking
entries with missing parent context, incompatible aliases, or only weak prose
evidence as ambiguous or out of scope. This conservative policy favors
precision over recall because an incorrect field merge can mislead all
downstream alignment and reasoning.

\subsection{Field-Centric Rule Extraction}
\label{sec:rules}

After field-space construction, \tool{} extracts normative requirements as
field-centric rules. The extraction prompt is constrained by the available field
space and by a fixed schema. The model must identify the constrained field,
condition, action, modality, source span, and any dependent fields. A
deterministic checker then validates that the field appears in the field space,
that the source span exists, and that the action belongs to a supported action
family.

\begin{table}[t]
\centering
\caption{Rule schema used by \tool.}
\label{tab:ruleschema}
\small
\begin{tabular}{lp{0.68\columnwidth}}
\toprule
\textbf{Attribute} & \textbf{Description} \\
\midrule
Field & Constrained protocol field or parameter \\
Condition & Message type, state, phase, or field predicate \\
Action & Presence, value, derivation, relation, or rejection \\
Modality & MUST, MUST NOT, SHOULD, MAY, or equivalent phrase \\
Evidence & Source section, sentence span, and referenced text \\
Dependencies & Other fields or states needed to evaluate the rule \\
\bottomrule
\end{tabular}
\end{table}

The action families are intentionally close to implementation behavior:
\emph{assign}, \emph{derive}, \emph{select}, \emph{require-present},
\emph{require-absent}, \emph{compare}, \emph{range-check},
\emph{state-update}, and \emph{reject}. This design encourages a rule to map to
observable code constructs such as assignments, conditionals, error returns, and
state transitions. When a sentence contains multiple constraints, \tool{}
atomizes it into multiple rules so that each rule can be aligned and tested
independently.

\textbf{Rule normalization.}
Before alignment, rules are normalized to make equivalent requirements compare
more consistently. Negated forms are rewritten into explicit actions, e.g.,
``MUST NOT accept $x$'' becomes an expected \emph{reject} action under condition
$x$. Cross-references are expanded when the referenced section contains a
field-space entry. Multiple modalities in the same sentence are split into
separate rules. The normalized rule remains linked to the original text so that
reviewers can audit whether the transformation preserved the intended meaning.

\textbf{Unsupported rules.}
Some normative statements are intentionally not converted into field-level
rules. Examples include broad performance guidance, implementation advice
without observable behavior, and cryptographic assumptions that require a proof
rather than code-level checking. \tool{} records these statements as
\emph{out-of-scope normative text}. Reporting out-of-scope text is useful: it
prevents silent loss of requirements and makes the method's coverage explicit.

\subsection{Syntax-Aware Rule-to-Code Alignment}
\label{sec:alignment}

The alignment stage maps a rule $r$ to implementation variables and code
contexts. It proceeds in three steps.

\textbf{File ranking.}
\tool{} uses an LLM-based ranker to score candidate C/C++ files by their
semantic relationship to the rule. The ranker sees the rule, field aliases,
action family, and file-level summaries. It cannot introduce files outside the
provided implementation. The top-ranked files are deduplicated and passed to
localization.

\textbf{Variable localization.}
For the constrained field $f$, \tool{} generates aliases by splitting field
tokens, normalizing case conventions, expanding common abbreviations, and
including parent-structure names when available. It then constructs a symbol
index using Tree-sitter~\cite{tree_sitter} with C/C++ grammars and
compiler-aware parsing where available. Candidate identifiers are ranked by
exact alias matches, token overlap, syntactic role, and proximity to
rule-action terms. Identifiers that are too generic, appear only as loop
variables, or lack a plausible syntactic relationship to the action are
down-ranked.

\textbf{Context retrieval.}
For each candidate variable $v$, \tool{} retrieves code snippets from three
sources: exact textual hits, symbol scopes containing the hits, and semantic
windows whose tokens overlap the rule. Each snippet $c$ receives an evidence
score:
\begin{equation}
\begin{aligned}
S_{\mathrm{evidence}}(c)
&=
5 \cdot \mathbb{I}_{\mathrm{exact}}(c)
 +0.9 \cdot |\mathcal{T}_c \cap \mathcal{T}_v|\\
&\quad
 +0.28 \cdot |\mathcal{T}_c \cap \mathcal{T}_r|
 +b(c).
\end{aligned}
\end{equation}
where $\mathbb{I}_{\mathrm{exact}}$ denotes exact variable or alias occurrence,
$\mathcal{T}_c$ is the code-token set, $\mathcal{T}_v$ is the variable-token set,
$\mathcal{T}_r$ is the rule-token set, and $b(c)$ rewards action-relevant usage
patterns such as initialization, validation, comparison, error handling, and
configuration reads. Snippets without exact variable evidence receive a score
cap to avoid weak semantic matches dominating direct code evidence.

The top snippets are then sent to an LLM semantic filter. The filter must choose
the snippets most relevant to the rule, state the relationship between the
snippet and the rule, and return a confidence score. If the model output is
malformed or empty, \tool{} falls back to the deterministic ranking.

The constants in Equation~(1) are intentionally simple rather than learned per
project. The exact-match bonus gives priority to snippets that mention the
localized field or one of its aliases; the variable-token term keeps renamed or
split representations in the candidate set; the rule-token term is weaker
because prose overlap alone often retrieves comments or unrelated helper
functions. We tune the three constants once on held-out standard sections and
keep them fixed across all evaluated protocols and implementations.

\textbf{Why not pure retrieval?}
Pure semantic retrieval often retrieves a function whose comments resemble the
standard but whose variables are unrelated to the constrained field. Pure exact
matching has the opposite problem: it misses implementations that rename fields
or distribute a field across parser state and connection state. \tool{} uses
syntax evidence as the meeting point. A snippet is strongest when it both
mentions a plausible representation of the field and uses that representation
in a role consistent with the rule action.

\subsection{Evidence-Driven Consistency Reasoning}
\label{sec:reasoning}

Once candidate contexts are available, \tool{} asks a reasoning agent to decide
whether the implementation satisfies the rule. The input contains the rule,
source provenance, selected snippets, call relationships discovered from the
symbol index, configuration conditions, and a compact trace of how the variable
was localized. The agent must produce a structured decision record with a
label, supporting evidence, assumptions, missing evidence, and a short defect
hypothesis when applicable.

\begin{algorithm}[t]
\small
\DontPrintSemicolon
\KwIn{Rule $r$, implementation index $I$, maximum rounds $K$}
\KwOut{Consistency decision record}
$C_0 \leftarrow$ AlignRuleToCode$(r,I)$\;
\For{$i \leftarrow 0$ \KwTo $K$}{
  $o_i \leftarrow$ Reason$(r,C_i)$\;
  \If{$o_i.label \neq \mathrm{Insufficient}$}{
    \Return $o_i$\;
  }
  $q_i \leftarrow$ MissingEvidenceQuery$(o_i)$\;
  $C_{i+1} \leftarrow C_i \cup$ RetrieveAdditionalContext$(q_i,I)$\;
  \If{$C_{i+1}=C_i$}{
    \Return $o_i$\;
  }
}
\Return $o_K$\;
\caption{Evidence-driven consistency reasoning.}
\label{alg:reasoning}
\end{algorithm}

Algorithm~\ref{alg:reasoning} shows the loop. The key design choice is to make
\textsc{Insufficient} a valid and useful outcome. If the current evidence lacks
a caller, macro, configuration option, or error path, the agent should not force
a binary judgment. Instead, it emits a missing-evidence query that triggers
another retrieval round. The loop stops when the agent reaches a determinate
label, no new evidence can be retrieved, or a maximum number of rounds is
reached.

\textbf{Evidence requests.}
Missing-evidence requests are typed. The current prototype supports requests
for callers, callees, macro definitions, configuration guards, field
assignments, validation checks, error constructors, state updates, and test
entry points. Typed requests make the loop less open-ended: a request for
``more code'' is rejected, while a request for ``the callers of
\texttt{parse\_transport\_parameter} that inspect its return value'' can be
served by the code indexer.

\textbf{Judgment policy.}
The reasoner may return \textsc{Inconsistent} only when it can point to both a
normative obligation and implementation behavior that contradicts the
obligation. Absence of evidence is not sufficient unless the search has covered
the relevant parser, validator, caller, and configuration paths specified by
the rule. This policy is conservative, but it is necessary for a system whose
reports are meant to be actionable for maintainers.

\subsection{Validation-Test Agent}
\label{sec:testagent}

Static reasoning can still be wrong: a missing branch may be enforced by a
caller, or an apparent check may be compiled out. \tool{} therefore validates
high-risk reports using generated tests. The test agent receives the rule, the
defect hypothesis, relevant snippets, build instructions, and any existing test
interfaces. It produces a test plan that attempts to construct one of two
counterexamples: an input that violates the standard but is accepted, or an
input that satisfies the standard but is rejected.

The agent then prepares the environment, selects build options, generates or
adapts a harness, executes the test, and collects logs. A report is confirmed
only when the observed behavior is stable and contradicts the field rule. If the
test cannot reach the relevant path, \tool{} records the unmet validation
conditions rather than upgrading the static hypothesis into a confirmed issue.
This conservative policy is important for reviewability: every reported
confirmed inconsistency is tied to either executable evidence or a clearly
marked manual reproduction plan.

\textbf{Oracle construction.}
The validation oracle is derived from the field rule, not from an existing
implementation. For a rejection rule, the oracle checks whether the target
returns the required error, closes the connection, or otherwise refuses the
input. For a presence or value rule, the oracle checks emitted messages,
configuration state, or parser outputs when those are observable. When the
expected behavior is not directly observable, \tool{} records the test as
inconclusive rather than inferring correctness from the absence of a crash.

\subsection{Implementation}
\label{sec:implementation}

We implemented \tool{} as a modular analysis pipeline. The current prototype
contains four components: a standard parser, a code indexer, an LLM orchestration
layer, and a validation-test runner. This section describes implementation
choices that affect reproducibility and engineering trade-offs.

\textbf{Standard processing.}
The standard parser ingests RFC-style text, markdown, or converted PDF text. It
segments documents by section, preserves sentence offsets, and detects
normative keywords using a conservative lexicon derived from RFC 2119-style
language. Field candidates are extracted from message diagrams, bullet lists,
tables, and definitional paragraphs. When a field candidate is ambiguous, the
parser retains all candidates and asks the extraction stage to resolve the
reference against local context.

\textbf{Code indexing.}
The code indexer supports C and C++ in the prototype. It uses Tree-sitter with
C/C++ grammars for a uniform symbol index and compiler-aware parsing where
available for richer scope information. Each indexed symbol records its file,
scope, declaration span, use spans, surrounding comments, and enclosing
function or type. The indexer also extracts a lightweight call graph from direct
calls and function-pointer-like patterns when they are syntactically visible.

\textbf{LLM orchestration.}
LLM calls are isolated behind typed prompts with JSON-like output schemas. The
pipeline uses low-temperature decoding for extraction, alignment, and
reasoning. Every LLM output passes through a deterministic validator that checks
schema validity, source-span existence, identifier presence, duplicate entries,
and label membership. Invalid outputs are retried with a compact repair prompt;
after the retry budget is exhausted, \tool{} falls back to deterministic
rankings or marks the sample as insufficient evidence.

\textbf{Prompt boundaries.}
The prompts are intentionally narrow. The extractor sees a section and the
field space, not the whole standard. The aligner sees a rule and ranked code
snippets, not the whole repository. The reasoner sees the evidence bundle and
must cite snippet identifiers rather than free-form code claims. This design
keeps token cost manageable and makes the output auditable.

\textbf{Validation runner.}
The validation runner currently supports three modes: extending an existing
unit test, creating a small harness around a parser or validation function, and
driving a command-line implementation with generated inputs. The runner records
compiler options, environment variables, command lines, exit status, stdout,
stderr, and relevant logs. If the target cannot be built or the path cannot be
reached, the sample is not counted as confirmed.

\textbf{Parameters.}
Unless otherwise stated, the prototype keeps the top 8 files after file
ranking, the top 20 variable candidates per rule, and the top 12 snippets per
reasoning round. The maximum number of evidence-completion rounds is three. The
validation agent receives at most two attempts per hypothesis before reporting
unmet validation conditions. We tune these parameters on held-out protocol
sections and keep them fixed for all evaluated implementations.

\textbf{Caching and determinism.}
To make repeated runs auditable, \tool{} stores every prompt, response,
validator decision, and retrieved code bundle under a content hash. When the
same rule and repository snapshot are analyzed again, deterministic stages are
reused and LLM stages can be replayed from cache. The cache is also useful for
manual inspection: a reviewer can open a report and see the exact evidence
available to the reasoner at each round.

\textbf{Failure handling.}
The implementation distinguishes recoverable and terminal failures. Recoverable
failures include malformed model output, missing optional build metadata, or a
snippet that cannot be parsed by one parser but can still be indexed by
Tree-sitter. Terminal failures include missing source files, standards that
cannot be segmented, or implementations that cannot expose any parse or
validation entry point. Terminal failures are reported as coverage gaps rather
than silently omitted from aggregate metrics.

\textbf{Manual review interface.}
Although \tool{} is designed for automation, the prototype emits a compact
review bundle for each candidate. The bundle contains the source text, the
normalized rule, ranked variables, selected snippets, reasoning transcript,
validation commands, and final label. This interface is important for building
ground truth: protocol experts and maintainers can correct a rule, mark an
implementation-specific exception, or confirm that a test demonstrates a real
deviation.
